A.~Eddington and F.~Dyson from Principe, and C.~Davidson and A.~Crommelin from
Sobral, Brazil measured the deflection of light coming from stars apparently
very close to the Sun during the total solar eclipse in 1919. The deflection
was found to agree with the theoretically predicted value of $1.75''$.

A light ray (or photon) which passes the Sun at a distance $d$ is deflected by
an angle
\[ \Delta\theta \propto \frac{4GM_\odot}{d c^2} \]

The present-day accuracy of the VLBI (Very Long Baseline Interferometry)
technique in the radio wavelength range is 0.1\,mas (milliarcsecond). Is it
possible to detect a deflection of radio photons from a quasar by (a) Jupiter,
(b) the Moon? Estimate the angle of deflection in both cases and mark ``YES''
or ``NO'' on the answer sheet.
